\documentclass[11pt]{article}

\usepackage[utf8]{inputenc}
\usepackage{color}
\usepackage{graphicx}
\usepackage{geometry}
\usepackage{tabularx}
\usepackage{hyperref}
\hypersetup{
    colorlinks,
    citecolor=black,
    filecolor=black,
    linkcolor=black,
    urlcolor=black
}

\renewcommand{\contentsname}{Inhaltsverzeichnis}

\begin{document}

\begin{titlepage}
	\begin{center}
		\vspace*{1cm}

		\Huge
		\textbf{Tourney}\\
		Spezifikation

		\vspace{0.5cm}
		\LARGE
		Softwarepraktikum\\
		\Large
		Wintersemester 2014/2015

		\vspace{1.5cm}

		\large
		\textbf{Jonas Auer (2860992)\\
				 Fabian Biester (2859084)\\
				 Jan Tagscherer (2893134)}

		\vfill

		\includegraphics[width=0.4\textwidth]{Logo.png}

		\vspace{1.5cm}

		\Large
		Universität Stuttgart\\
		14.11.2014
	\end{center}
\end{titlepage}

\newpage

\tableofcontents
\newpage

\section{Einleitung}

\subsection{Zweck der Spezifikation}

Dieses Dokument dient dem Zweck, sowohl funktionale als auch qualitative Anforderungen überprüfbar, vollständig und systematisch festzuhalten. Damit ermöglicht es dem Kunden und den Betreuern sicherzustellen, dass die richtige Funktionalität auf korrekte Art und Weise umgesetzt wird. Um dies zu gewährleisten muss die Spezifikation als wichtiger Leitfaden während der Entwicklung und für alle weiteren erstellten Artefakte beachtet werden. Zudem müssen sowohl die Spezifikation als auch weitere betreffende Dokumente aktuell gehalten werden.

\subsection{Leserkreis}

Der Leserkreis dieses Dokuments setzt sich aus den folgenden Personengruppen zusammen:
\begin{itemize}
	\item Das Projektteam:
	\begin{itemize}
		\item Die Entwickler
		\item Die Gutachter während des Reviews
		\item Die Betreuer
	\end{itemize}
	\item Der Auftraggeber
	\item Mit der Wartung und Weiterentwicklung betraute Software-Entwickler
\end{itemize}

\subsection{Einsatzbereich und Ziele}

In diesem Projekt wird die Organisationssoftware \textbf{Tourney} umgesetzt. Diese setzt sich zum Ziel, die Planung und Durchführung von Spielturnieren deutlich zu vereinfachen.

Zuvor fand die Organisation dieser Turniere weitestgehend manuell statt, indem Anmeldungen und Ergebnisse auf Papier oder per Tabellenkalkulation aufgezeichnet wurden. Dies führte zu starken Verzögerungen und deutlich mehr Arbeit im Rahmen eines solchen Turniers.

Mit dem Einsatz von \textbf{Tourney} entfällt unnötiger Mehraufwand, indem Events und deren Turniere modular verwaltet werden. Dazu gehören sowohl das Führen von Datenbe-ständen wie Voranmeldungen, Anmeldungen und Turnierergebnissen als auch das automatische Auswerten von Spielresultaten anhand bearbeitbarer Spielregeln und die Verteilung einiger Daten über Turnierordner oder weitere administrative Arbeitsplätze mit anschließendem Zusammenfügen.

Die Software wird im Rahmen des Softwarepraktikums im Wintersemester 2014/2015 an der Universität Stuttgart und unter Auftrag des Heidelberger Spieleverlags entwickelt.

\subsection{Fachbegriffe und Abkürzungen}

Relevante Fachbegriffe und deren Abkürzungen sind in der Spezifikation \textit{kursiv} markiert und werden im Begriffslexikon geklärt, das diesem Dokument anhängt.

\subsection{Übersicht und Aufbau des Dokuments}

Um die zu entwickelnde Software vollständig zu spezifizieren wird das Dokument in die nachfolgenden Kapitel gegliedert:
\begin{itemize}
	\item[] \textbf{Kapitel 2} gibt einen groben Überblick über grundlegende Funktionen, über die die Software verfügen soll, in welchem Umfeld und von welchen Nutzern sie eingesetzt wird und welche Einschränkungen und Annahmen bei der Entwicklung eingegangen werden.
	\item[] \textbf{Kapitel 3} beschreibt die konkreten Anforderungen, die vom Auftraggeber an die Software gestellt werden. Dabei werden zwischen funktionalen und qualitativen Anforderungen unterschieden.
	\item[] \textbf{Kapitel 4} charakterisiert Nutzer und für die Software relevante Personengruppen.
	\item[] \textbf{Kapitel 5} enthält Skizzen als Prototyp der geplanten Benutzeroberfläche.
	\item[] \textbf{Kapitel 6} umfasst grundlegende Anwendungsfälle, die bei der Benutzung der Software auftreten können und die die Interaktion der Nutzer mit dem Programm umreißen.
	\item[] \textbf{Kapitel 7} umschließt ein Begriffslexikon, in dem alle wichtigen Fachbegriffe und Abkürzungen geklärt werden.
\end{itemize}

\section{Allgemeine Beschreibung}

\subsection{Einbettung}

Bisher wurden keine Softwarelösungen zur Organisation von Spielturnieren eingesetzt. Deshalb muss \textbf{Tourney} in keine bestehenden Systeme eingebettet werden.

\subsection{Einschränkungen bei der Entwicklung}

Aufgrund der geforderten Plattformunabhängigkeit wird die Software in der Programmiersprache Java entwickelt. Dabei können die Versionen 1.7 oder 1.8 verwendet werden.

Zusätzlich können die folgenden Bibliotheken verwendet werden:
\begin{itemize}
	\item Grafische Benutzeroberfläche mit \textbf{Swing} oder \textbf{JavaFX}
	\item Datenbankeinbindung durch \textbf{H2} in Verbindung mit \textbf{JDBC}
	\item Erstellen von PDF-Dokumenten mit \textbf{iText}
\end{itemize}

\subsection{Grundlegende Produktfunktionen}

Die Software soll grundsätzlich die folgende Funktionalität haben, die in Kapitel 3 ausführ-licher beschrieben wird:
\begin{itemize}
	\item Modulare Erstellung von Regelmodulen für konkrete Spielturniere
	\item Planung von Events und Festlegen von stattfindenden Turnieren
	\item Verwalten von Spielerlisten in Hinsicht auf Voranmeldung, Anmeldung bei Anwesenheit und Bezahlstatus
	\item Durchführung und Auswertung von Turnieren anhand spezifizierten Regeln und An-zeigen von Paarungen über eine Projektion
	\item Zusammenführen und Ausgeben der Turnierresultate am Administrationsrechner
	\item Verteilen von relevanten Datenbanken zum kollaborativen Arbeiten während Anmeldungen und an die Turnierordner mit jeweils anschließendem Zusammenfügen der Daten
\end{itemize}

\subsection{Benutzermerkmale}

Die Nutzer von \textbf{Tourney} sollten zur erfolgreichen Bedienung keine besonderen Kenntnisse besitzen müssen. Insbesondere sollte die Oberfläche so selbsterklärend sein, dass dies auch ohne Lesen des Handbuchs möglich ist. Die Benutzer besitzen hierbei für gewöhnlich zwar Fachkenntnis in Bezug auf die Spielturniere, allerdings sollte kein besonderes technisches Wissen vorausgesetzt werden. Vielmehr sollte die Software so intuitiv zu bedienen sein, dass für keinen Nutzer ein Hindernis entsteht.

Desweiteren kann nicht davon ausgegangen werden, dass sowohl die Arbeitsplätze von Administratoren als auch Turnierordnern ausreichend abgesichert sind. Deshalb sollten relevante Daten über Passwörter zusätzlich gesichert sein.

Zuletzt sollte versehentlichem Datenverlust durch eine übergreifende Rückgängig-Funktion vorgebeugt werden.

\subsection{Annahmen und Abhängigkeiten}

Um die Software erfolgreich entwickeln und betreiben zu können muss von den folgenden externen Einflussfaktoren ausgegangen werden:
\begin{itemize}
	\item Die erforderliche Hardware zur Ausführung des Programms und Massenspeichergeräte zur Datenübertragung müssen zur Laufzeit bereitgestellt werden.
	\item Die Turnierregeln dürfen sich nicht so stark ändern, dass die modulare Regelerstellung nicht mehr relevant ist. In diesem Fall kann eine korrekte Ergebnisauswertung nicht mehr garantiert werden.
	\item Aufgrund des engen zeitlichen Rahmens dürfen in späten Projektphasen keine fundamentalen Änderungen an den Anforderungen mehr erfolgen.
\end{itemize}

\section{Spezifische Anforderungen}

\subsection{Funktionale Anforderungen}

\subsubsection{Mengengerüst}

\subsubsection{Leistungsanforderungen}

\subsubsection{Externe Schnittstellen}

\subsubsection{...}

\subsection{Qualitätsanforderungen}

\subsubsection{Bedienbarkeit}

\subsubsection{Erweiterbarkeit}

\subsubsection{Sicherheit}

\subsubsection{Zuverlässigkeit}

\subsubsection{Portabilität}

\subsubsection{...}

\section{Benutzerprofile}

\subsection{Administrator}

Der Administrator hat übergreifenden Zugriff über die erstellten Datensätze. Er kann Events erstellen, Turniere hinzufügen und verwaltet Spieleranmeldungen. Diese Daten kann er teilweise auch im Nachhinein bei Bedarf abändern. Zudem kann der Administrator Auszüge aus den erstellten Datensätzen an weitere Benutzer weitergeben, einerseits um Anmeldungen an mehreren Arbeitsplätzen vorzunehmen, andererseits um Spielerlisten und Turniereigenschaften an die Turnierordner zu übergeben. Nach diesen Vorgängen werden die erhobenen Daten wieder beim Administrator zusammengefügt.

Beim Administrator handelt es sich um den Organisator des Events. Auch bei dieser Person werden keine besonderen technischen Kenntnisse vorausgesetzt.

\subsection{Turnierordner}

Turnierordner haben die Aufgabe, tatsächlich anwesende Spieler bei den jeweiligen einzelnen Turnieren zu erfassen und dann die Spiele durchzuführen. Dabei werden von der Software Paarungen vorgeschlagen und nach dem Spielende die Resultate eingetragen. Diese Daten lassen sich auch später noch manuell abändern. Zudem kann der Turnierordner Spieler disqualifizieren. Für diese Vorgänge werden die vom Administrator übertragenen Daten verwendet.

Ein Turnierordner ist zumeist eine Person mit guten Kenntnissen der Spielregeln und sollte die Verantwortung über die Datensätze übernehmen können.

\subsection{Spieler}

Die Spieler sind die größte Nutzergruppe. Sie führen keine direkte Interaktion mit der Software durch, sondern sollen sich über Ausgaben und Projektionen über Paarungen und Spielstände informieren können. Dazu müssen diese leicht verständlich und weithin lesbar sein.

\section{Benutzeroberfläche}

\subsection{...}

\section{Anwendungsfälle}

\subsection{Übersicht über die Anwendungsfälle}

\subsection{Anwendungsfälle des Administrators}

\subsubsection{Event erstellen}

\begin{tabularx}{\textwidth}{| p{0.2\textwidth} | p{0.7415\textwidth} |}
	\hline
	\textbf{Ziel} & ... \\
	\hline
	\textbf{Akteure} &
		\begin{enumerate}
			\item ...
			\item ...
		\end{enumerate}
	\\
	\hline
	\textbf{Vorbedingung} & ... \\
	\hline
	\textbf{Normalablauf} &
		\begin{enumerate}
			\item ...
			\item ...
		\end{enumerate}
	\\
	\hline
	\textbf{Nachbedingung} & ... \\
	\hline
	\textbf{Sonderfälle} &
		\begin{enumerate}
			\item ...
			\item ...
		\end{enumerate}
	\\
	\hline
	\textbf{Nachbedingung im Sonderfall} & ... \\
	\hline
\end{tabularx}

\subsection{Anwendungsfälle der Turnierordner}

\section{Anhang}

\subsection{Begriffslexikon}

\subsection{Versionshistorie}

\begin{itemize}
	\item Version 1.0 (05.11.2014)
	\begin{itemize}
		\item Grundstruktur der Spezifikation erstellt
	\end{itemize}
	\item Version 1.1 (10.11.2014)
	\begin{itemize}
		\item Kapitel 1 und 2 hinzugefügt
	\end{itemize}
	\item Version 1.2 (11.11.2014)
	\begin{itemize}	
		\item Kapitel 4 hinzugefügt
	\end{itemize}
\end{itemize}

\end{document}
